Selenium nanoparticles (SeNPs) have attracted widespread attention due to their favorable biocompatibility and diverse bioactivities; however, their cytotoxic responses are co-regulated by particle size, concentration, surface functionalization, and material chemical states, rendering systematic evaluation via conventional experimental assays labor-intensive and low-throughput. To address the challenge of predicting selenium nanoparticle cytotoxicity, this study establishes a two-tiered research framework: "benchmark dataset model comparison + dedicated dataset enhanced-feature modeling." First, leveraging a publicly available multi-type nanoparticle dataset curated by Masarkar et al. \cite{masarkar2024}, a systematic performance comparison was conducted across Random Forest, XGBoost, Attention-Based Random Forest (ABRF), Leaf Attention-Based Random Forest (LARF), TabTransformer, and distilled TabTransformer architectures. Subsequently, based on a curated domain-specific dataset encompassing 383 validated selenium nanoparticle records, systematic feature statistics, preprocessing pipelines, and predictive models were developed.

Empirical results demonstrate that on large-scale heterogeneous tabular data, TabTransformer achieved the highest $R^2$ (0.768) in regression tasks, whereas tree-based ensembles exhibited superior stability in classification tasks, with XGBoost yielding an area under the receiver operating characteristic curve (AUC) of 0.945. On the dedicated selenium nanoparticle dataset—extracted from 43 peer-reviewed publications spanning 2008--2022—the incorporation of material-chemistry-related enhanced features markedly improved the fitting capacity of tree-based models compared to the baseline feature configuration. Specifically, XGBoost yielded optimal performance under the enhanced feature set, exhibiting a three-fold cross-validation mean absolute error (MAE) of 7.20% and an $R^2$ of 0.957; Random Forest likewise achieved high predictive precision (MAE of 7.63%, $R^2$ of 0.915). Automated feature selection revealed that the combination of IC50, elemental Se content, long-term stability, and Se 3d binding energy constituted the optimal subset of enhanced variables. Nonetheless, given the intrinsic codependency between IC50 and cytotoxicity endpoints, its predictive contribution is more appropriately designated as auxiliary validation evidence rather than an independent predictor situated at the same hierarchical tier as a priori physicochemical descriptors.

This investigation demonstrates that relying exclusively on generic physicochemical descriptors is insufficient for nanotoxicity prediction within single-material systems; instead, enhanced features capturing intrinsic chemical attributes offer essential complementary value. The data curation, feature characterization, and modeling workflows established herein provide a methodological foundation for developing selenium nanoparticle toxicity prediction platforms, prioritizing laboratory screening candidates, and guiding subsequent external validation studies, while findings incorporating IC50 require cautious interpretation tailored to realistic prospective prediction scenarios.

% 纳米硒（SeNPs）因具有良好的生物相容性和多重生物活性而受到广泛关注，但其毒性效应受粒径、浓度、表面修饰和材料化学状态等因素共同影响，传统实验方法难以快速完成系统评估。本文围绕纳米硒细胞毒性预测问题，构建了“基准数据集模型比较 + 专用数据集增强特征建模”的研究框架。首先，在来源于 Masarkar 等人的多类型纳米颗粒公开数据集 \cite{masarkar2024} 上，对随机森林、XGBoost、ABRF、LARF、TabTransformer 和蒸馏 TabTransformer 进行了系统比较；随后基于我们整理的 383 条硒纳米颗粒有效记录，进一步开展专用数据集特征统计、数据预处理分析与模型构建。

% 结果表明，在大样本异质表格数据上，TabTransformer 在回归任务中取得最高 $R^2$（0.768），而树模型在分类任务中表现更稳定，其中 XGBoost 的 AUC 为 0.945。在硒纳米颗粒专用数据集上，数据来源覆盖 43 篇文献、时间跨度为 2008--2022 年，引入材料化学相关增强信息后，树模型的拟合能力较基础特征设置有所改善，其中 XGBoost 在增强特征设置下取得较优结果，三折交叉验证 MAE 为 7.20\%，$R^2$ 为 0.957；随机森林在增强特征下也达到较高精度（MAE 7.63\%，$R^2$ 0.915）。进一步的自动筛选表明，IC50、元素 Se 含量、长期稳定性和结合能 Se3d 构成表现较优的增强变量组合；但考虑到 IC50 与细胞毒性终点之间存在天然关联，其贡献更适合视为辅助验证信息，而非与先验理化特征完全同层级的独立预测变量。

% 本研究说明，面向单一材料体系的毒性预测，仅依赖通用理化参数是不够的，能够刻画材料本征化学属性的增强特征具有重要补充价值。所构建的数据整理、特征统计和建模流程可为纳米硒毒性预测平台开发、实验优先级筛选及后续外部验证研究提供方法基础，而含 IC50 的结果仍需结合真实预测场景进一步审慎解释。
